\subsubsection{Reconfiguration energy law}
The simulated windows give a smooth law (Fig.~\ref{fig:X7}): for a fixed lateral displacement the energy above isolated cruise falls as the window lengthens (lower lateral speed, smaller heading change, smaller airspeed excess) while the wake benefit lost grows with the time spent outside the slot. For the \SI{369.6}{m} split the net overhead per follower is \novFast\,Wh at \ntFast\,s, \novMid\,Wh at \ntMid\,s and \novSlow\,Wh at \ntSlow\,s; the \SI{60}{m} join-up and break-up cost under \SI{0.1}{Wh} each. The maximum bank angle in these manoeuvres is $8$--$10^\circ$ for the long splits and $13$--$21^\circ$ for the short joins (peak lateral speed \SIrange{8}{12}{m/s}), within the \SI{35}{\degree} limit. Merge and split are close to mirror images of each other, differing mainly in where the wake benefit is regained or lost.
\begin{figure}[H]\centering\includegraphics[width=\textwidth]{X7_reconfiguration_law}
\caption{Follower energy above isolated cruise versus (a) window duration and (b) peak lateral speed; (c) maximum bank angle. Splits (circles) and merges (squares) for $Y=\SI{60}{m}$ and \SI{369.6}{m}.}\label{fig:X7}\end{figure}
